\chapter{Worried about your memory}
\index{Worried about your memory}
\section*{Definition and Overview}
Memory lapses are common and do not necessarily mean dementia. Concern may reflect normal ageing, stress, poor sleep, depression, anxiety, medication effects, or a cognitive disorder. Progressive decline that affects everyday function deserves assessment.
\section*{Clinical Features}
Concerning changes include repeatedly forgetting recent conversations, getting lost in familiar places, difficulty managing medicines or money, language problems, personality change, or loss of previously safe skills. Sudden confusion is different from gradual memory decline and may be delirium or stroke.
\section*{Differential Diagnosis}
Consider sleep deprivation or sleep apnoea, depression, anxiety, grief, medication or alcohol effects, thyroid disease, vitamin B12 deficiency, infection, delirium, stroke, epilepsy, head injury, mild cognitive impairment, and dementia.
\section*{When to Seek Medical Care}
Arrange a routine appointment for persistent or progressive concerns, or if family members notice a change. Seek emergency care for sudden confusion, new weakness, facial droop, speech difficulty, severe headache, seizure, fever, or reduced consciousness.
\section*{Diagnosis and Investigations}
Assessment includes the person's history, an informant history, medicines and substance use, mood and sleep, physical and neurological examination, and cognitive testing. Blood tests and brain imaging are selected according to findings. No single test diagnoses dementia in every person.
\section*{Management}
Treat reversible contributors, review medicines, manage vascular risks, improve sleep, support hearing and vision, and address depression or anxiety. If cognitive impairment is confirmed, referral to a memory or specialist service can clarify cause and provide planning and support.
\section*{Complications}
Unrecognised cognitive impairment can increase risks involving driving, falls, medicines, finances, work, and exploitation. Anxiety about memory can itself worsen concentration.
\section*{Prognosis and Outcomes}
Many causes of memory difficulty are treatable or improve when contributing factors are addressed. Neurodegenerative conditions are progressive but benefit from early diagnosis, support, and advance planning.
\section*{Prevention and Self-Care}
Sleep regularly, remain physically and socially active, manage blood pressure and diabetes, avoid smoking and harmful alcohol use, use hearing aids when prescribed, and keep a written list of symptoms and medicines for assessment.
\section*{Sources and Further Reading}
\begin{itemize}
    \item NHS memory loss: \url{https://www.nhs.uk/conditions/memory-loss-amnesia/}
    \item National Institute on Aging: \url{https://www.nia.nih.gov/health/memory-loss}
\end{itemize}
